\begin{theorem}[Bayes--InfoNCE 的 Laplace--Rényi 表示与离散矩谱有限展开]\label{thm:fold-infonce-laplace-renyi-moment-expansion}
\leanverified{paper\_fold\_infonce\_laplace\_renyi\_moment\_expansion}
沿用定义 \ref{def:fold-pushforward-weight-wm} 的推前分布 \(w_m\)，并沿用定理 \ref{thm:fold-infonce-entropy-degeneration-second-order} 的采样机制：取 \(K\ge 2\)，先采样 \(X\sim w_m\)，再令
\[
J\mid X\sim \mathrm{Binomial}\bigl(K-1,\ w_m(X)\bigr),
\qquad
L_m^{K,\star}=\EE\bigl[\log(1+J)\bigr].
\]
定义 Rényi 幂和（碰撞矩谱）
\[
M_q(m):=\sum_{x\in X_m} w_m(x)^q=\EE_{X\sim w_m}\bigl[w_m(X)^{q-1}\bigr]\qquad(q\ge 1).
\]
则成立绝对收敛的积分表示
\[
L_m^{K,\star}
=
\int_{0}^{\infty}\frac{e^{-t}}{t}
\Bigl(1-\EE\bigl[(1-(1-e^{-t})w_m(X))^{K-1}\bigr]\Bigr)\,\dd t.
\]
并且 \(L_m^{K,\star}\) 具有有限阶矩谱展开
\[
L_m^{K,\star}
=
\sum_{j=1}^{K-1}(-1)^{j-1}\binom{K-1}{j}\,\alpha_j\,M_{j+1}(m),
\qquad
\alpha_j:=\int_{0}^{\infty}\frac{e^{-t}(1-e^{-t})^j}{t}\,\dd t\in(0,\infty).
\]
\end{theorem}
\begin{proof}
对任意整数 \(n\ge 0\) 有恒等式
\[
\log(1+n)=\int_{0}^{\infty}\frac{e^{-t}-e^{-(1+n)t}}{t}\,\dd t
=\int_{0}^{\infty}\frac{e^{-t}}{t}\bigl(1-e^{-nt}\bigr)\,\dd t.
\]
当 \(n\in\{0,1,\dots,K-1\}\) 时，上式右端由
\(
e^{-t}\min\{n,1/t\}
\)
一致控制，故可对随机变量 \(J\) 交换期望与积分，得到
\[
L_m^{K,\star}
=
\int_{0}^{\infty}\frac{e^{-t}}{t}\Bigl(1-\EE[e^{-tJ}]\Bigr)\,\dd t.
\]
对条件分布 \(J\mid X\sim\mathrm{Binomial}(K-1,w_m(X))\)，其 Laplace 变换为
\[
\EE\bigl[e^{-tJ}\mid X\bigr]
=\bigl(1-w_m(X)+w_m(X)e^{-t}\bigr)^{K-1}
=\bigl(1-(1-e^{-t})w_m(X)\bigr)^{K-1}.
\]
代回即得积分表示。
\par
再以二项式定理展开
\[
1-(1-cW)^{K-1}=\sum_{j=1}^{K-1}(-1)^{j-1}\binom{K-1}{j}c^j W^j,
\qquad c=1-e^{-t},\ \ W=w_m(X),
\]
并注意 \(\EE[W^j]=\sum_x w_m(x)^{j+1}=M_{j+1}(m)\)。
将该展开代入积分表示并逐项积分（有限和）得到所述矩谱展开。证毕。
\end{proof}

\begin{corollary}[对数系数闭式：Bayes--InfoNCE 作为“对数--矩谱”有限线性泛函]\label{cor:fold-infonce-alpha-log-closed-form}
在定理 \ref{thm:fold-infonce-laplace-renyi-moment-expansion} 的记号下，对每个 \(j\ge 1\) 有完全显式的闭式
\[
\alpha_j
=
-\sum_{k=0}^{j}(-1)^k\binom{j}{k}\log(k+1)
=
-\sum_{k=1}^{j+1}(-1)^{k-1}\binom{j}{k-1}\log k.
\]
因此对任意整数 \(K\ge 2\)，\(L_m^{K,\star}\) 落在由 \(\{\log 2,\log 3,\dots,\log K\}\) 张成的 \(\QQ\)-线性空间中；
其系数由有限阶矩谱 \(\{M_2(m),\dots,M_K(m)\}\) 代数确定。
\end{corollary}
\begin{proof}
由定义
\(
\alpha_j=\int_{0}^{\infty}\frac{e^{-t}(1-e^{-t})^j}{t}\dd t
\)。
对 integrand 作二项式展开：
\[
e^{-t}(1-e^{-t})^j=\sum_{k=0}^{j}(-1)^k\binom{j}{k}e^{-(k+1)t}.
\]
并注意系数和为 \(\sum_{k=0}^{j}(-1)^k\binom{j}{k}=(1-1)^j=0\)。
因此可用 Frullani 型恒等式：若 \(a_k>0\) 且 \(\sum_k c_k=0\)，则
\[
\int_{0}^{\infty}\frac{\sum_k c_k e^{-a_k t}}{t}\,\dd t=-\sum_k c_k\log a_k,
\]
从而得到
\(
\alpha_j=-\sum_{k=0}^{j}(-1)^k\binom{j}{k}\log(k+1)
\)。
第二个求和式仅为换元 \(k\mapsto k-1\)。
最后一句由定理 \ref{thm:fold-infonce-laplace-renyi-moment-expansion} 的有限展开直接推出。证毕。
\end{proof}
\leanverified{paper\_fold\_infonce\_alpha\_log\_closed\_form}

\begin{proposition}[Bayes--InfoNCE 的幂和展开与通用差分系数]\label{prop:fold-infonce-bayes-infonce-power-sum-expansion}
固定分辨率 \(m\) 并令 \(p_r:=M_r(m)=\sum_{x\in X_m}w_m(x)^r\)（\(r\ge 1\)）。
对每个整数 \(r\ge 2\) 定义差分系数
\[
b_r:=\sum_{j=0}^{r-1}(-1)^{r-1-j}\binom{r-1}{j}\log(j+1).
\]
则对任意整数 \(K\ge 2\) 有幂和展开
\[
L_m^{K,\star}=\sum_{r=2}^{K}\binom{K-1}{r-1}\,b_r\,p_r.
\]
并且 \(b_r\) 具有积分表示
\[
b_r=(-1)^r\int_{0}^{\infty}\frac{e^{-t}(1-e^{-t})^{r-1}}{t}\,\dd t,
\]
从而 \(b_r\neq 0\) 且 \(\mathrm{sign}(b_r)=(-1)^r\)（\(r\ge 2\)）。
\end{proposition}
\begin{proof}
由定理 \ref{thm:fold-infonce-laplace-renyi-moment-expansion} 的有限展开，令 \(r=j+1\) 得
\[
L_m^{K,\star}
=\sum_{r=2}^{K}(-1)^{r-2}\binom{K-1}{r-1}\,\alpha_{r-1}\,p_r.
\]
定义 \(b_r:=(-1)^r\alpha_{r-1}\)（注意 \(r-2\equiv r\pmod 2\)），即得所示幂和展开。
由推论 \ref{cor:fold-infonce-alpha-log-closed-form} 可知
\(
\alpha_{r-1}=-\sum_{j=0}^{r-1}(-1)^j\binom{r-1}{j}\log(j+1)
\)，从而得到 \(b_r\) 的差分闭式。
积分表示由 \(\alpha_{r-1}\) 的积分定义与 \(b_r=(-1)^r\alpha_{r-1}\) 得到。
由于被积函数严格正，\(\alpha_{r-1}>0\) 且 \(\mathrm{sign}(b_r)=(-1)^r\)。证毕。
\end{proof}
\leanverified{paper\_fold\_infonce\_bayes\_infonce\_power\_sum\_expansion}

\begin{corollary}[损失塔对幂和的三角反演]\label{cor:fold-infonce-loss-tower-triangular-inversion-power-sums}
在命题 \ref{prop:fold-infonce-bayes-infonce-power-sum-expansion} 的记号下，对每个整数 \(q\ge 2\) 有递推反演
\[
p_q
=\frac{1}{b_q}\left(
L_m^{q,\star}-\sum_{k=2}^{q-1}\binom{q-1}{k-1}\,b_k\,p_k
\right).
\]
因此对任意 \(Q\ge 2\)，有限前缀 \(\{L_m^{2,\star},\dots,L_m^{Q,\star}\}\) 唯一确定 \(\{p_2,\dots,p_Q\}\)。
\end{corollary}
\begin{proof}
由命题 \ref{prop:fold-infonce-bayes-infonce-power-sum-expansion} 在 \(K=q\) 的特例，
\(
L_m^{q,\star}=\sum_{k=2}^{q}\binom{q-1}{k-1}b_k p_k
\)。
注意对角系数 \(\binom{q-1}{q-1}b_q=b_q\neq 0\)，移项即得递推式。证毕。
\end{proof}

\begin{theorem}[Bayes--InfoNCE 最优损失谱的矩谱反演与纤维谱完备性]\label{thm:fold-infonce-loss-spectrum-identifiability}
\leanverified{paper\_fold\_infonce\_loss\_spectrum\_identifiability}
固定分辨率 \(m\ge 1\)，并记 \(N:=|X_m|\)。
对每个整数 \(K\ge 2\)，令 \(L_m^{K,\star}\) 表示定理 \ref{thm:fold-infonce-entropy-degeneration-second-order} 的 Bayes 最优 InfoNCE 最小损失。
则对一切 \(2\le K\le N\) 存在显式系数 \(\beta_{K,q}\in\RR\)（\(2\le q\le K\)）使得
\begin{equation}\label{eq:fold-infonce-loss-spectrum-triangular}
L_m^{K,\star}=\sum_{q=2}^{K}\beta_{K,q}\,M_q(m),
\qquad
M_q(m):=\sum_{x\in X_m} w_m(x)^q,
\end{equation}
并且该线性系统在未知量 \(\{M_q(m)\}_{2\le q\le N}\) 上为下三角且主对角元非零。
因此最优损失谱 \(\{L_m^{K,\star}\}_{K=2}^{N}\) 唯一确定
\begin{enumerate}
\normalfont
  \item 碰撞矩有限段 \(\{M_q(m)\}_{q=2}^{N}\)；
  \item 纤维多重度多重集合 \(\{d_m(x):x\in X_m\}\) 与推出权重多重集合 \(\{w_m(x):x\in X_m\}\)（均按重数计）。
\end{enumerate}
更精确地，对每个 \(2\le q\le N\)，\(M_q(m)\) 仅依赖于子谱 \(\{L_m^{K,\star}\}_{2\le K\le q}\)。
\end{theorem}
\begin{proof}
由定理 \ref{thm:fold-infonce-laplace-renyi-moment-expansion}，
\[
L_m^{K,\star}
=
\sum_{j=1}^{K-1}(-1)^{j-1}\binom{K-1}{j}\,\alpha_j\,M_{j+1}(m),
\qquad
\alpha_j\in(0,\infty).
\]
令 \(q=j+1\)，并定义
\[
\beta_{K,q}:=(-1)^{q-2}\binom{K-1}{q-1}\alpha_{q-1}\qquad(2\le q\le K),
\]
则得到 \eqref{eq:fold-infonce-loss-spectrum-triangular}。
由于 \(\beta_{K,K}=(-1)^{K-2}\alpha_{K-1}\neq 0\)，系统为下三角且可逐步递推反演，从而 \(\{L_m^{K,\star}\}_{K=2}^{N}\) 唯一确定 \(\{M_q(m)\}_{q=2}^{N}\)。
\par
由定义 \ref{def:fold-pushforward-weight-wm}，\(w_m(x)=d_m(x)/2^m\)，故
\[
S_q(m):=\sum_{x\in X_m}d_m(x)^q=2^{qm}M_q(m)\qquad(q\ge 1),
\]
并且 \(S_1(m)=\sum_x d_m(x)=|\Omega_m|=2^m\) 由 \(m\) 已知。
因此由 \(\{M_q(m)\}_{2\le q\le N}\) 可恢复 \(\{S_q(m)\}_{1\le q\le N}\)。
由定理 \ref{thm:op-algebra-watatani-index-spectrum-identifiable-from-collision-moments}，该有限段碰撞矩唯一确定多重集合 \(\{d_m(x)\}_{x\in X_m}\)，从而亦确定 \(\{w_m(x)\}_{x\in X_m}\)。证毕。
\end{proof}

\begin{corollary}[损失塔的 Newton--Prony 完备性：权重谱、碰撞矩与直方图]\label{cor:fold-infonce-loss-tower-newton-prony-completeness}
在定理 \ref{thm:fold-infonce-loss-spectrum-identifiability} 的设定下，记 \(p_r:=M_r(m)=\sum_{x\in X_m}w_m(x)^r\) 且 \(p_1=1\)。
则最优损失谱 \(\{L_m^{K,\star}\}_{K=2}^{N}\) 代数确定如下对象：
\begin{enumerate}
  \item \textnormal{(Newton)} 单项首一多项式
  \[
  Q_m(z):=\prod_{x\in X_m}(z-w_m(x)),
  \]
  其根（按重数计）给出多重集合 \(\{w_m(x):x\in X_m\}\)，并由 \(d_m(x)=2^m w_m(x)\) 恢复 \(\{d_m(x)\}\)；
  \item \textnormal{(Collision moments)} 对任意整数 \(q\ge 1\)，碰撞矩
  \[
  S_q(m):=\sum_{x\in X_m}d_m(x)^q
  =2^{qm}p_q
  \]
  可由损失谱有限前缀 \(\{L_m^{K,\star}\}_{K=2}^{q}\) 递推计算得到；
  \item \textnormal{(Histogram)} 若 \(\{d_m(x)\}\) 的不同取值为 \(D_1<\cdots<D_r\) 且
  \(n_i:=\#\{x\in X_m:\ d_m(x)=D_i\}\)，
  则有限矩列 \(\{S_k(m)\}_{k=0}^{2r-1}\) 在代数意义下唯一恢复 \(\{(D_i,n_i)\}_{i=1}^{r}\)
  （Prony--Hankel 反演；参见命题 \ref{prop:conclusion-tqft-genus-linear-recurrence-recovery} 的同型机制）。
\end{enumerate}
\end{corollary}
\begin{proof}
由定理 \ref{thm:fold-infonce-loss-spectrum-identifiability} 可恢复 \(\{p_2,\dots,p_N\}\)，再由 Newton 恒等式以 \(\{p_1,\dots,p_N\}\) 唯一恢复 \(Q_m\)，得 (1)。
由 \(d_m(x)=2^m w_m(x)\) 得 (2) 的缩放恒等式，而 \(p_q\) 首次出现在 \(K=q\) 的下三角关系中，故只需有限前缀即可递推求出。
最后 (3) 为有限指数和矩量反演的标准结论。证毕。
\end{proof}
\leanverified{paper\_fold\_infonce\_loss\_tower\_newton\_prony\_completeness}

\begin{corollary}[由 InfoNCE 损失谱恢复指数元、配分函数与预言机容量]\label{cor:fold-infonce-loss-spectrum-determines-index-tqft-capacity}
\leanverified{paper\_fold\_infonce\_loss\_spectrum\_determines\_index\_tqft\_capacity}
在定理 \ref{thm:fold-infonce-loss-spectrum-identifiability} 的设定下，给定最优损失谱 \(\{L_m^{K,\star}\}_{K=2}^{N}\) 进而唯一确定：
\begin{enumerate}
\normalfont
  \item Watatani 指标元 \(\mathrm{Ind}_{\mathrm W}(E_m)=\sum_{x\in X_m}d_m(x)\widetilde P_x\)（定理 \ref{thm:op-algebra-fold-watatani-index-equals-multiplicity-field}）及其迹矩
  \(
  \tau_m(\mathrm{Ind}_{\mathrm W}(E_m)^q)=2^{-m}\sum_x d_m(x)^{q+1}
  \)
  （推论 \ref{cor:op-algebra-fold-watatani-index-moments}）；
  \item 任意整数 \(g\ge 0\) 的二维配分函数
  \(
  Z_m(\Sigma_g)=\sum_{x\in X_m}d_m(x)^{2-2g}
  \)
  （定理 \ref{thm:conclusion-fold-symtft-partition-function-collision-moments}）；
  \item 预言机容量曲线与成功率曲线
  \(
  \mathcal{C}_m(B)=\sum_x \min(d_m(x),2^B)
  \)
  与 \({\rm Succ}_m(B)=\mathcal{C}_m(B)/2^m\)
  （定理 \ref{thm:fold-oracle-capacity-closed-form}），并可将 \(\sum_x w_m(x)^q\) 代入定理 \ref{thm:oracle-renyi-nonasymptotic-error-bound} 得到任意 \(q>1\) 的非渐近失败证书。
  \item 令
  \[
  D_m:=\max_{x\in X_m}d_m(x),
  \qquad
  N_m^\star:=\min\{N\in\NN_{\ge 1}:\ d(N)\ge D_m\},
  \]
  其中 \(d(\cdot)\) 为约数函数（定义 \ref{def:conclusion-truncated-prime-register}--定理 \ref{thm:conclusion-divisor-lattice-optimal-bit-scaling} 的约数格证书库口径）。
  则最优损失谱唯一确定 \(D_m\)，从而唯一确定最小 prime-register 状态预算 \(N_m^\star\) 及其比特预算 \(\log_2 N_m^\star\)；
  并且当 \(m\to\infty\) 时满足尺度律
  \[
  \log_2 N_m^\star
  =\left(\frac{\log_2\varphi}{2}+o(1)\right)m\log_2 m.
  \]
\end{enumerate}
\par
此外，把 \eqref{eq:fold-infonce-loss-spectrum-triangular} 与推论 \ref{cor:op-algebra-fold-watatani-index-moments} 代回可得显式“矩变换”恒等式：对任意 \(2\le K\le N\)，
\[
L_m^{K,\star}
=\sum_{q=2}^{K}\beta_{K,q}\,2^{-(q-1)m}\,\tau_m\!\left(\mathrm{Ind}_{\mathrm W}(E_m)^{q-1}\right).
\]
\end{corollary}
\begin{proof}
由定理 \ref{thm:fold-infonce-loss-spectrum-identifiability} 已可恢复多重集合 \(\{d_m(x)\}\)，而上述三类对象均为 \(\{d_m(x)\}\) 的函子化表达或幂和表达；逐条代回所引定理/命题即得。证毕。
\end{proof}

\leanverified{paper\_fold\_infonce\_triangular\_inversion\_collision\_spectrum}
\begin{corollary}[等价形式：差分算子表示与前差分系数]\label{cor:fold-infonce-triangular-inversion-collision-spectrum}
沿用命题 \ref{prop:fold-infonce-bayes-infonce-power-sum-expansion} 的记号，并令
\[
f(j):=\log(1+j),\qquad (\Delta f)(j):=f(j+1)-f(j)
\quad(j\in\ZZ_{\ge 0}).
\]
则对任意整数 \(n\ge 1\)（即取 \(K=n+1\)），有严格三角表示
\[
L_m^{n+1,\star}
=\sum_{t=1}^{n}\binom{n}{t}\,(\Delta^t f)(0)\,M_{t+1}(m),
\qquad
(\Delta^t f)(0)=\sum_{j=0}^{t}(-1)^{t-j}\binom{t}{j}\log(1+j)=b_{t+1}.
\]
并且最高阶系数非退化：\((\Delta^{n}f)(0)=(-1)^{n+1}\alpha_n\neq 0\)。
因此得到显式递归反演：对每个 \(n\ge 1\)，
\[
M_{n+1}(m)
=\frac{1}{(\Delta^{n}f)(0)}
\left(
L_m^{n+1,\star}-\sum_{t=1}^{n-1}\binom{n}{t}\,(\Delta^{t}f)(0)\,M_{t+1}(m)
\right).
\]
从而有限序列 \(\{L_m^{2,\star},\dots,L_m^{Q+1,\star}\}\) 唯一确定 \(\{M_2(m),\dots,M_{Q+1}(m)\}\)。
\end{corollary}
\begin{proof}
由定理 \ref{thm:fold-infonce-laplace-renyi-moment-expansion}，
\[
L_m^{n+1,\star}
=\sum_{j=1}^{n}(-1)^{j-1}\binom{n}{j}\,\alpha_j\,M_{j+1}(m).
\]
另一方面由推论 \ref{cor:fold-infonce-alpha-log-closed-form} 与 \((\Delta^j f)(0)=\sum_{k=0}^{j}(-1)^{j-k}\binom{j}{k}\log(1+k)\) 得
\(
(-1)^{j-1}\alpha_j=(\Delta^j f)(0)
\)，代回即得三角表示。
又由定理 \ref{thm:fold-infonce-laplace-renyi-moment-expansion} 的积分定义知 \(\alpha_n>0\)，故 \((\Delta^n f)(0)=(-1)^{n+1}\alpha_n\neq 0\)，从而可按 \(n\) 递归解出 \(M_{n+1}(m)\)。
证毕。
\end{proof}

\begin{theorem}[学习相变边界的不可判定性：停机问题嵌入 Bayes--InfoNCE 速率]\label{thm:fold-infonce-rate-undecidable}
固定 \(\alpha\in(0,1)\)，并令 \(K_m:=2^{\alpha m}\)。
考虑判定问题：输入为一个一致可计算的映射族 \(f_m:\Omega_m=\{0,1\}^m\to Y_m\)（\(m\ge 1\)），其诱导推前分布记为 \(w_m^{(f)}\)，并令 \(L_{m}^{K,\star}(f)\) 表示以 \(w_m^{(f)}\) 代替 \(w_m\) 后、按定理 \ref{thm:fold-infonce-laplace-renyi-moment-expansion} 所定义的 Bayes 最优 InfoNCE 损失。
判定是否成立
\[
\lim_{m\to\infty}\frac{1}{m}L_{m}^{K_m,\star}(f)=0.
\]
则该判定问题不可判定；更强地，该问题的真值集合既非递归可枚举，亦非其补集递归可枚举。
\end{theorem}
\begin{proof}
给定任意图灵机 \(M\) 与输入 \(x\)，对每个 \(m\) 先模拟 \(M(x)\) 运行 \(m\) 步。
若在 \(m\) 步内停机，则令 \(f_m\) 为恒等映射 \(f_m(\omega)=\omega\)（取 \(Y_m=\Omega_m\)）；
若在 \(m\) 步内不停机，则令 \(f_m\) 为常值映射（取 \(Y_m=\{\ast\}\) 且 \(f_m(\omega)=\ast\)）。
该构造对每个 \(m\) 仅需有限步模拟，故 \(m\mapsto f_m\) 一致可计算。
\par
若 \(M(x)\) 永不停机，则对一切 \(m\)，推前分布退化为 \(w_m^{(f)}(\ast)=1\)，从而 \(J=K_m-1\) 几乎处处，故
\[
L_{m}^{K_m,\star}(f)=\log K_m=\alpha m\log 2,
\qquad
\lim_{m\to\infty}\frac{1}{m}L_{m}^{K_m,\star}(f)=\alpha\log 2>0.
\]
若 \(M(x)\) 在某个有限步数 \(T\) 内停机，则对所有 \(m\ge T\)，\(f_m\) 恒等，推前分布为均匀分布，故 \(w_m^{(f)}(X)=2^{-m}\) 几乎处处且
\[
J\sim\mathrm{Binomial}(K_m-1,2^{-m}),
\qquad
\EE[J]=(K_m-1)2^{-m}=2^{(\alpha-1)m}\to 0.
\]
由 \(0\le \log(1+J)\le J\) 得
\(
0\le L_{m}^{K_m,\star}(f)\le \EE[J]\to 0
\)，从而
\(
\lim_{m\to\infty}m^{-1}L_{m}^{K_m,\star}(f)=0
\)。
因此上述极限判定与 \(M(x)\) 是否停机等价，故不可判定。
最后一条“既非 r.e.\ 亦非 co-r.e.” 由停机集合及其补集均非递归可枚举的经典结论导出。证毕。
\end{proof}

% Auto-split to keep each TeX source < 600 lines.

\begin{proposition}[Bayes--InfoNCE 学习相变边界泛函的统一可逼近性不可计算]\label{prop:fold-infonce-boundary-functional-uncomputable}
\leanverified{paper\_fold\_infonce\_boundary\_functional\_uncomputable}
令 \(\mathcal F\) 为所有一致可计算的映射族 \(f=(f_m)_{m\ge 1}\) 的集合，
并令 \(L_m^{K,\star}(f)\) 表示以 \(f\) 诱导的推前分布 \(w_m^{(f)}\) 代替 \(w_m\) 后的 Bayes 最优 InfoNCE 损失（定理 \ref{thm:fold-infonce-laplace-renyi-moment-expansion}）。
对每个 \(f\in\mathcal F\) 定义边界泛函
\[
\mathfrak B(f)
:=
\sup\Bigl\{\alpha\in(0,1):\ \lim_{m\to\infty}\frac{1}{m}L_m^{2^{\alpha m},\star}(f)=0\Bigr\}\in[0,1].
\]
则不存在图灵机 \(\mathsf A\)，能够对任意输入 \((f,\epsilon)\in\mathcal F\times\QQ_{>0}\) 在有限步内输出有理数 \(r\) 使得
\[
|\mathfrak B(f)-r|<\epsilon.
\]
特别地，不存在对全体 \(f\in\mathcal F\) 同时有效的统一有限步数值近似程序，能够以任意预设精度逼近该学习相变边界。
\end{proposition}

\begin{proof}
反证。
若存在上述 \(\mathsf A\)，则取 \(\epsilon=\tfrac14\)。
对定理 \ref{thm:fold-infonce-rate-undecidable} 的停机嵌入构造：给定任意 \((M,x)\)，得到一致可计算族 \(f\) 满足
若 \(M(x)\) 停机，则 \(f_m\) 对充分大 \(m\) 为恒等映射，从而 \(\lim_{m\to\infty}m^{-1}L_m^{2^{\alpha m},\star}(f)=0\) 对一切 \(\alpha\in(0,1)\) 成立并推出 \(\mathfrak B(f)=1\)；
若 \(M(x)\) 不停机，则 \(f_m\) 恒为常值映射，从而 \(\lim_{m\to\infty}m^{-1}L_m^{2^{\alpha m},\star}(f)=\alpha\log 2>0\) 并推出 \(\mathfrak B(f)=0\)。
\par
因此 \(\mathsf A(f,\tfrac14)\) 的输出 \(r\) 满足 \(r>\tfrac12\) 当且仅当 \(M(x)\) 停机，从而给出停机问题的判定算法，矛盾。证毕。
\end{proof}

